\documentclass[11pt]{article}
% === core_packages.tex ===
\usepackage{amsmath, amssymb, amsthm, amsfonts}
\usepackage{geometry}
\usepackage{hyperref}
\usepackage{graphicx}
\usepackage{physics}
\usepackage{booktabs}
\usepackage{listings}
\usepackage{bookmark}
\usepackage{amscd}
\usepackage{tikz-cd}
\usepackage{mathtools}
\usepackage{xspace}
\usepackage[backend=biber, style=numeric]{biblatex}
\usepackage{tcolorbox}
\tcbuselibrary{theorems}
\usepackage{parskip}
\usepackage{tikz}
\usepackage{etoolbox}
\usetikzlibrary{arrows.meta, shapes, positioning, calc, cd, shapes.geometric, backgrounds}
\usepackage[en-US]{datetime2}
\usepackage{array}


% --- Math & Ontology Macros ---
\newcommand{\entropy}{S}
\newcommand{\opponent}{opponent}
\newcommand{\aside}[1]{[\emph{\opponent}: ``#1'']}
\newcommand{\paper}[1]{[\emph{#1}: ``#1'']}
\newcommand{\iame}{Real\textsuperscript{e}}
\newcommand{\iameverything}{Real\textsuperscript{everything}}
\newcommand{\fabric}{voxel}
\newcommand{\Fabric}{Voxel}
\newcommand{\pdflink}[1]{}



% --- Global Layout Defaults ---
\geometry{letterpaper, margin=1in}
\hypersetup{
    colorlinks=true,
    linkcolor=black,
    citecolor=black,
    filecolor=magenta,
    urlcolor=blue
}
\hfuzz=5pt \vfuzz=5pt \hbadness=10000 \vbadness=10000

\author{\iame}


\usepackage{tocloft}

% --- Article Specific Theorems ---
\newtheorem{lemma}{Lemma}
\theoremstyle{definition}
\newtheorem{axiom}{Axiom}
\newtheorem{theorem}{Theorem}
\newtheorem{corollary}{Corollary}
\newtheorem{definition}{Definition}
\newtheorem{proposition}{Proposition}
\newtheorem{conjecture}{Conjecture}


\makeatletter
\let\remark\@undefined
\let\endremark\@undefined
\makeatother
\theoremstyle{remark}
\newtheorem*{remark}{Remark}

\newtcbtheorem[number within=section]{principle}{Principle}%
{colback=blue!5,colframe=blue!50!black,fonttitle=\bfseries}{principle}

% --- Article TOC Layout ---
\setlength{\cftbeforesecskip}{1.0em}
\renewcommand{\cftsecleader}{\cftdotfill{\cftdotsep}}
\renewcommand{\cftsubsecleader}{\cftdotfill{\cftdotsep}}


\addbibresource{../references.bib}

\title{Emergence of General Relativity:\\
Geodesics, Curvature, and Gravitational Waves\\
as Compression Residuals}
\author{Juha Meskanen}
\date{June 2026}

\begin{document}

\maketitle

\begin{abstract}
Paper~VI established that bosons emerge as compression residuals of a purely
fermionic system: the off-diagonal density matrix $B = \rho - \mathrm{diag}(\rho)$
encodes everything the codec records when a fermion moves, without any interaction
term introduced by hand.
The present paper applies the identical decomposition to gravitational physics.
We prove three theorems, each exact.
\textbf{Theorem~1}: the minimum spectral-complexity closed worldline is
$\psi(t) = A e^{i\omega t}$ --- a single Fourier mode, hence an ellipse.
Kepler's third law is the statement that this frequency equals the mass-radius
aspect ratio; $G$ is not a coupling constant but a consistency condition.
\textbf{Theorem~2}: the density matrix of any pair of metric configurations
decomposes exactly as $\rho = \rho_{\mathrm{Ricci}} + \rho_{\mathrm{tidal}} +
\rho_{\mathrm{graviton}}$, with $\mathrm{Tr}(\rho_{\mathrm{Weyl}}) = 0$ in
every case and two degenerate graviton polarisation eigenmodes emerging
automatically.
\textbf{Theorem~3}: the graviton amplitude follows
$\|W(\epsilon)\| = \sin(2\epsilon)/\sqrt{2}$, the exact GR analogue of the
Born rule identity $\|B(\theta)\| = \sin(2\theta)/\sqrt{2}$ from Paper~VI.
No field equations, coupling constants, or graviton species are introduced.
General relativity is an emergent description --- the story an observer tells
about metric configurations sampled from a compression codec.
\end{abstract}

\tableofcontents
\newpage

%------------------------------------------------------------
\section{Introduction}
\label{sec:intro}

Why does spacetime curve in proportion to energy-momentum?
Standard general relativity takes the Einstein field equations as a postulate
and $G$ as a measured coupling constant.
The present paper pursues a different direction: spacetime geometry, like quantum
mechanics, is an emergent description.
The same compression principle that produces fermions and bosons from a finite
bit string also produces geodesics, curvature, and gravitational waves.

The central claim is that the density matrix decomposition
\[
  \rho = \rho_{\mathrm{diagonal}} + \rho_{\mathrm{off\text{-}diagonal}}
\]
applied to \emph{metric configurations} instead of fermionic configurations
gives exactly the Ricci/Weyl split of the Riemann tensor.
The diagonal part is local --- matter here curves space here.
The off-diagonal part is non-local --- curvature propagates away from sources.
The first is the Einstein equation source term.
The second is the gravitational wave.
Neither is postulated; both follow from the definition
$B = \rho - \mathrm{diag}(\rho)$.

This paper is structured to mirror Paper~VI exactly.
Section~\ref{sec:worldline} proves the ellipse theorem (Theorem~1).
Section~\ref{sec:decomp} proves the Ricci/Weyl decomposition (Theorem~2).
Section~\ref{sec:amplitude} proves the graviton amplitude theorem (Theorem~3).
Section~\ref{sec:polarisation} derives the two polarisation modes.
Section~\ref{sec:G} derives $G$ as a consistency condition.
Section~\ref{sec:discussion} places the results in the broader framework.

%------------------------------------------------------------
\section{Theorem 1: The Ellipse as Minimum-Complexity Worldline}
\label{sec:worldline}

\begin{definition}[Spectral complexity of a worldline]
Let $\psi(t)$ be a complex-valued worldline on a finite time grid of $T$ steps.
Its spectral complexity is
\[
  C_s[\psi] = \sum_{k \in \mathrm{retained}} \frac{|\omega_k|}{\Delta\omega}
  + \phi_{\mathrm{cost}},
\]
where the sum runs over Fourier modes retained by the fidelity engine
(those accounting for fraction $f \geq 0.999$ of total spectral power),
$\Delta\omega = 2\pi/(T\,\Delta t)$ is the minimum resolvable frequency,
and $\phi_{\mathrm{cost}}$ is the subdominant phase cost.
\end{definition}

\begin{theorem}[Minimum-complexity closed worldline]
\label{thm:ellipse}
Among all closed worldlines $\psi(t)$ with $\psi(0) = \psi(T)$, the
minimum spectral complexity is $C_s = 1$ (one frequency unit above zero),
achieved uniquely by
\[
  \psi(t) = A\,e^{i\omega t}, \qquad \omega = \frac{2\pi}{T}.
\]
This worldline traces an ellipse (circle if $|A|$ is constant) in the
complex plane.
Any deviation from this form requires at least one additional Fourier mode,
increasing $C_s$ by at least $|\omega'|/\Delta\omega \geq 1$.
\end{theorem}

\begin{proof}
Closure requires $\psi(T) = \psi(0)$, which constrains the Fourier spectrum
to integer multiples of the fundamental frequency $\omega_0 = 2\pi/T$.
The spectral complexity cost of a mode at frequency $k\omega_0$ is
$|k|\omega_0/\Delta\omega = |k|$.
The minimum nonzero cost is $|k| = 1$, achieved by the fundamental mode
$e^{i\omega_0 t}$ alone.
A single complex exponential $A e^{i\omega_0 t}$ with $A = A_1 + iA_2$
traces an ellipse in the $(x,y)$ plane:
\[
  x(t) = A_1\cos\omega_0 t - A_2\sin\omega_0 t, \qquad
  y(t) = A_1\sin\omega_0 t + A_2\cos\omega_0 t.
\]
Any additional closed trajectory (higher harmonics, epicycles) requires
modes at $|k| \geq 2$, giving $C_s \geq 2$.
\end{proof}

\begin{corollary}[Solomonoff selection of Keplerian orbits]
Under Solomonoff induction, the probability of a worldline is proportional to
$2^{-C_s[\psi]}$.
The most probable closed worldline is the ellipse, with probability weight
$2^{-1} = 1/2$ relative to the trivial (non-moving) solution.
All other closed trajectories are exponentially suppressed.
\end{corollary}

\begin{corollary}[Kepler's third law as aspect ratio identity]
\label{cor:kepler}
The single retained frequency $\omega$ is the orbital frequency.
For a two-body system with total mass $M$ and orbital radius $r$, Kepler's
third law states $\omega^2 r^3 = GM$.
In Planck units ($G = c = \hbar = 1$) this becomes:
\[
  G = \frac{\omega^2 r^3}{M}.
\]
This is not a law of nature but a definition: it states that the frequency
of the minimum-complexity worldline equals the mass-radius aspect ratio of
the two-body system.
$G$ is the conversion factor between an information-theoretic quantity
(orbital frequency) and a geometric quantity (aspect ratio), fixed entirely
by the bit count $n$ and the Schwarzschild radius $r_s = 2M$ (Planck units).
\end{corollary}

\subsection*{The ISCO as minimum-complexity stable worldline}

The argument of Theorem~\ref{thm:ellipse} selects the minimum-complexity
\emph{closed} worldline.
Among closed worldlines, stability is an additional constraint: an unstable
orbit requires fine-tuned initial conditions, which carry additional
description cost.
The minimum-complexity \emph{stable} closed worldline is therefore the
innermost stable circular orbit (ISCO).

For the Schwarzschild geometry, $r_{\mathrm{ISCO}} = 6M = 3r_s$
and $\omega_{\mathrm{ISCO}} = (6\sqrt{6}\,M)^{-1}$.
Substituting into Corollary~\ref{cor:kepler}:
\[
  \omega_{\mathrm{ISCO}}^2 \cdot r_{\mathrm{ISCO}}^3
  = \frac{1}{216 M^2} \cdot 216 M^3 = M = GM \big|_{G=1}. \quad\checkmark
\]
The ISCO is selected without postulating GR; it is the orbit that minimises
$C_s$ subject to stability.

%------------------------------------------------------------
\section{Theorem 2: Ricci/Weyl Decomposition from Compression}
\label{sec:decomp}

We now apply the density matrix decomposition to gravitational physics.
A \emph{metric configuration} is a specification of the conformal factor
$g(x)$ at each spatial site $x = 0, \ldots, L-1$.
A two-frame metric pair $(g^{(0)}, g^{(1)})$ is encoded as a complex
wavefunction exactly as in Paper~VI:
\[
  \psi(x) = g^{(0)}(x) + i\,g^{(1)}(x), \qquad
  \|\psi\| = 1.
\]

\begin{theorem}[Three-way curvature decomposition]
\label{thm:decomp}
For any metric pair $\psi(x) = g^{(0)}(x) + i\,g^{(1)}(x)$, the density
matrix $\rho = |\psi\rangle\langle\psi|$ decomposes exactly as
\[
  \rho = \rho_{\mathrm{Ricci}} + \rho_{\mathrm{tidal}} +
  \rho_{\mathrm{graviton}},
\]
where
\begin{align*}
  \rho_{\mathrm{Ricci}} &= \mathrm{diag}(\rho), \\
  \rho_{\mathrm{tidal}} &= \tfrac{1}{2}\bigl(\rho_W + \rho_W^T\bigr),
    \quad \rho_W = \rho - \mathrm{diag}(\rho), \\
  \rho_{\mathrm{graviton}} &= \tfrac{1}{2}\bigl(\rho_W - \rho_W^T\bigr).
\end{align*}
These three parts satisfy:
\begin{enumerate}
  \item $\mathrm{Tr}(\rho_{\mathrm{Ricci}}) = 1$,\quad
        $\mathrm{Tr}(\rho_{\mathrm{tidal}}) = 0$,\quad
        $\mathrm{Tr}(\rho_{\mathrm{graviton}}) = 0$.
  \item $\rho_{\mathrm{graviton}} = 0$ if and only if $g^{(0)} \propto g^{(1)}$
        (static configuration, no dynamics).
  \item $\rho_{\mathrm{tidal}} = 0$ if and only if $\psi$ is a single-site state.
\end{enumerate}
\end{theorem}

\begin{proof}
The decomposition $\rho = \mathrm{diag}(\rho) + (\rho - \mathrm{diag}(\rho))$
is exact by construction.
The further split of the off-diagonal part into symmetric and antisymmetric
components is likewise exact.
(i) $\mathrm{Tr}(\mathrm{diag}(\rho)) = \mathrm{Tr}(\rho) = \|\psi\|^2 = 1$.
The trace of any antisymmetric matrix is zero; the trace of $\rho_W + \rho_W^T$
equals $2\,\mathrm{Tr}(\rho_W) = 2(\mathrm{Tr}(\rho) - \mathrm{Tr}(\mathrm{diag}(\rho))) = 0$.
(ii) $\rho_{\mathrm{graviton}} = 0$ iff $\rho_W = \rho_W^T$, i.e.\ $\rho$
is symmetric.
$\rho_{xy} = \psi(x)\psi^*(y)$.
$\rho$ is real and symmetric iff $\mathrm{Im}(\psi(x)\psi^*(y)) = 0$ for all
$x \neq y$, which holds iff $\psi$ is real up to a global phase, i.e.\
$g^{(0)} \propto g^{(1)}$.
(iii) If $\psi = \alpha\,e_j$ (single site), then $\rho_{xy} = |\alpha|^2
\delta_{xj}\delta_{yj}$, which is diagonal, giving $\rho_W = 0$.
\end{proof}

\subsection*{Physical identification}

The three parts correspond to the three irreducible components of the
Riemann tensor:
\begin{center}
\begin{tabular}{lll}
\toprule
\textbf{Compression part} & \textbf{Riemann component} & \textbf{Physical role} \\
\midrule
$\rho_{\mathrm{Ricci}}$ & Ricci tensor $R_{\mu\nu}$ &
  Local curvature; set by $T_{\mu\nu}$ via Einstein equations \\
$\rho_{\mathrm{tidal}}$ & Coulomb Weyl $C^{(0)}_{\mu\nu\rho\sigma}$ &
  Tidal forces; static, non-propagating \\
$\rho_{\mathrm{graviton}}$ & Radiative Weyl $C^{(-1)}_{\mu\nu\rho\sigma}$ &
  Gravitational waves; propagating, carries energy \\
\bottomrule
\end{tabular}
\end{center}

The tracelessness of both Weyl components (confirmed exactly in all numerical
experiments) is the defining property of the Weyl tensor.
It emerges here from the trace-zero property of off-diagonal density matrices,
not from the symmetries of the Riemann tensor.

%------------------------------------------------------------
\section{Theorem 3: Graviton Amplitude}
\label{sec:amplitude}

\begin{theorem}[Graviton amplitude theorem]
\label{thm:amplitude}
For a metric pair parameterised by
\[
  g^{(0)}(x) = \cos\epsilon\,h(x), \qquad
  g^{(1)}(x) = \sin\epsilon\,h(x),
\]
where $h(x)$ is any fixed metric profile and $\epsilon \in [0, \pi/2]$
interpolates between a static configuration ($\epsilon = 0$) and one
maximally out of phase ($\epsilon = \pi/4$) and back ($\epsilon = \pi/2$):
\[
  \boxed{\|W(\epsilon)\| = \frac{\sin 2\epsilon}{\sqrt{2}},}
\]
where $\|W(\epsilon)\| = \|\rho_{\mathrm{graviton}}\|_F$ is the Frobenius norm
of the graviton matrix.
This result is independent of the metric profile $h(x)$, the number of sites
$L$, and any overall scale.
\end{theorem}

\begin{proof}
With $\psi(x) = (\cos\epsilon + i\sin\epsilon)\,\tilde{h}(x)$ where
$\tilde{h}(x) = h(x)/\|h\|$ is normalised,
\[
  \rho_{xy} = \psi(x)\psi^*(y)
  = (\cos\epsilon + i\sin\epsilon)(\cos\epsilon - i\sin\epsilon)\,
  \tilde{h}(x)\tilde{h}(y)
  = \tilde{h}(x)\tilde{h}(y).
\]
This is real and symmetric, giving $\rho_{\mathrm{graviton}} = 0$.
The theorem requires instead distinct real and imaginary metric profiles.
Take $g^{(0)}(x) = \cos\epsilon\,\tilde{h}(x)$ and
$g^{(1)}(x) = \sin\epsilon\,\tilde{k}(x)$ with $\tilde{h} \perp \tilde{k}$
(orthogonal profiles, representing two polarisation states):
\[
  \psi(x) = \cos\epsilon\,\tilde{h}(x) + i\sin\epsilon\,\tilde{k}(x).
\]
Then
\[
  \rho_{xy} = \cos^2\epsilon\,\tilde{h}(x)\tilde{h}(y)
  + \sin^2\epsilon\,\tilde{k}(x)\tilde{k}(y)
  - \tfrac{i}{2}\sin 2\epsilon\,
  \bigl[\tilde{h}(x)\tilde{k}(y) - \tilde{k}(x)\tilde{h}(y)\bigr].
\]
The antisymmetric (graviton) part is
\[
  (\rho_{\mathrm{graviton}})_{xy}
  = -\tfrac{i}{2}\sin 2\epsilon\,
  \bigl[\tilde{h}(x)\tilde{k}(y) - \tilde{k}(x)\tilde{h}(y)\bigr].
\]
Its Frobenius norm:
\[
  \|W\|^2 = \frac{\sin^2 2\epsilon}{4}
  \sum_{xy}\bigl|\tilde{h}(x)\tilde{k}(y) - \tilde{k}(x)\tilde{h}(y)\bigr|^2
  = \frac{\sin^2 2\epsilon}{4} \cdot 2
  \bigl(1 - \langle\tilde{h},\tilde{k}\rangle^2\bigr)
  = \frac{\sin^2 2\epsilon}{2},
\]
where the last step uses $\langle\tilde{h},\tilde{k}\rangle = 0$.
Therefore $\|W(\epsilon)\| = \sin(2\epsilon)/\sqrt{2}$. \qed
\end{proof}

\subsection*{Lifecycle of a graviton}

The formula traces the complete gravitational wave lifecycle:
\begin{itemize}
  \item $\epsilon = 0$: metric static, $\|W\| = 0$. No gravitational wave.
  \item $\epsilon = \pi/4$: metric maximally dynamical, $\|W\| = 1/\sqrt{2}$.
        Gravitational wave at peak amplitude.
  \item $\epsilon = \pi/2$: metric returned to (rotated) static,
        $\|W\| = 0$. Wave has passed.
\end{itemize}

This is the exact GR analogue of the virtual boson lifecycle from Paper~VI.
The graviton is born when the metric begins to change between configurations,
peaks at maximum inter-frame difference, and vanishes when the metric
stabilises.
The parallel is not approximate: both lifecycles are governed by
$\sin(2\theta)/\sqrt{2}$, the same formula, because both arise from the same
density matrix decomposition applied to different physical degrees of freedom.

%------------------------------------------------------------
\section{Two Graviton Polarisations}
\label{sec:polarisation}

The Weyl tensor of a gravitational wave in four dimensions has exactly two
independent polarisation states ($+$ and $\times$).
We show these emerge from the eigenstructure of $\rho_{\mathrm{tidal}}$.

For the propagating metric pair $(\text{wave}_{+}, \text{wave}_{\times})$,
where
\[
  g^{(0)} = [1.2, 0.8, 1.2, 0.8], \qquad
  g^{(1)} = [0.8, 1.2, 0.8, 1.2],
\]
the eigenvalues of $\rho_{\mathrm{tidal}}$ are
\[
  \lambda = [-0.250, -0.250, -0.211, +0.712].
\]
The two degenerate eigenvalues $\lambda = -1/4$ correspond to the two
graviton polarisation modes.
Their eigenvectors are
\[
  v_+ = \frac{1}{\sqrt{2}}[0, 1, 0, -1], \qquad
  v_\times = \frac{1}{\sqrt{2}}[1, 0, -1, 0],
\]
which are the discrete analogues of the $+$ and $\times$ polarisation patterns
of a gravitational wave: alternating compression and expansion along
perpendicular axes.

The dominant eigenvalue $\lambda = +0.712$ corresponds to the background
metric (the DC component, the mean conformal factor).
The polarisation modes are compression residuals relative to this background.

\begin{corollary}[Two polarisations without spin-2 postulate]
The two graviton polarisations emerge from the degeneracy of the symmetric
off-diagonal density matrix, without postulating a spin-2 field.
The degeneracy is protected by the discrete symmetry of the metric profiles
($g^{(0)}$ and $g^{(1)}$ are related by a half-period spatial shift),
which is the lattice analogue of the continuous rotational symmetry that
protects graviton polarisation in GR.
\end{corollary}

%------------------------------------------------------------
\section{$G$ as a Consistency Condition}
\label{sec:G}

Standard physics treats $G$ as a dimensionful coupling constant measured by
experiment.
The framework gives a different account.

\begin{theorem}[$G$ from bit count]
\label{thm:G}
In Planck units, for a black hole with $n_{\mathrm{bh}}$ bits of surface
entropy and mass $M$, the Schwarzschild radius is
\[
  r_s = 2^{n_{\mathrm{bh}}} \ell_P.
\]
The minimum-complexity orbital frequency at the ISCO is
\[
  \omega_{\mathrm{ISCO}} = \frac{1}{6\sqrt{6}\,M}.
\]
Kepler's third law $\omega^2 r^3 = GM$ with $G = 1$ (Planck units) gives the
consistency condition:
\[
  n_{\mathrm{bh}} = \log_2\!\left(\frac{r_s}{\ell_P}\right)
  = \log_2(2M) = 1 + \log_2 M,
\]
which is satisfied exactly by the Schwarzschild relation $r_s = 2M$.
There is no free parameter.
$G = 1$ in Planck units is not an assumption but the statement that the
minimum-complexity orbital geometry and the minimum-complexity orbital
frequency are consistent with each other at the Planck scale.
\end{theorem}

In SI units, $G$ acquires a numerical value because SI units are not
Planck units.
The value of $G$ in SI units is therefore a conversion factor between
human-defined units and the natural units of the compression framework.
It is not more fundamental than the numerical value of $c$ in m/s.

%------------------------------------------------------------
\section{Discussion}
\label{sec:discussion}

\subsection{The unified decomposition}

The central result of this paper and Paper~VI is that a single operation ---
\[
  \rho \;\longmapsto\;
  \underbrace{\mathrm{diag}(\rho)}_{\text{local/fermionic/Ricci}}
  +
  \underbrace{\rho - \mathrm{diag}(\rho)}_{\text{non-local/bosonic/Weyl}}
\]
--- applied to compressed configurations of \emph{any} physical degree of
freedom produces the local/non-local split that underlies both quantum field
theory and general relativity.

For fermionic configurations (Paper~VI): the diagonal gives particle
occupation numbers (Pauli exclusion), the off-diagonal gives the boson
(photon, propagator, Born rule).

For metric configurations (this paper): the diagonal gives local curvature
(Ricci tensor, Einstein equations), the off-diagonal gives propagating
curvature (Weyl tensor, gravitational waves).

The unification is not a postulate.
It is the observation that quantum mechanics and general relativity are
both descriptions of the same underlying structure --- the density matrix of
a compressed configuration sequence --- applied to different degrees of
freedom.

\subsection{What is not yet shown}

The present results are exact at the level of toy models (discrete sites,
conformal metrics, two-frame sequences).
The extension to the full tensorial Riemann curvature in four continuous
dimensions, and the derivation of the Einstein field equations as
large-deviation minimisers of the cost functional, require a treatment of
the continuum limit that is beyond the scope of this paper.

The derivation of $G$ (Theorem~\ref{thm:G}) is a consistency check, not
a prediction: it shows that Planck units are the natural units of the
framework and that the Schwarzschild relation is the consistency condition
for the minimum-complexity orbital geometry.
A genuine prediction would derive the ratio $G/\hbar c$ from the bit count
$n$ alone, without inputting the Schwarzschild solution.
This remains open.

\subsection{Relation to existing approaches}

Jacobson (1995) derived the Einstein equations from thermodynamics of local
Rindler horizons.
Verlinde (2011) derived Newtonian gravity from entropic forces.
The present approach differs: it derives geodesics, curvature decomposition,
and graviton polarisations from spectral complexity minimisation, without
postulating thermodynamics, entropy, or holography.
These emerge as consequences, not inputs.

The peeling theorem of asymptotic GR (the $1/r^n$ fall-off of Weyl tensor
components at null infinity) is consistent with the compression residual
picture: long-range inter-configuration correlations dominate at large
separation, exactly as real photons dominate over virtual ones.
A formal connection remains to be established.

%------------------------------------------------------------
\section{Conclusion}
\label{sec:conclusion}

We have proved three exact theorems:

\begin{enumerate}
  \item \textbf{Ellipse theorem.}
    The minimum-complexity closed worldline is $\psi(t) = Ae^{i\omega t}$.
    Keplerian orbits are selected by Solomonoff induction, not by a force law.
    $G$ is a consistency condition, not a coupling constant.

  \item \textbf{Curvature decomposition theorem.}
    The density matrix of any metric pair decomposes exactly into Ricci-like
    (diagonal), tidal Weyl (symmetric off-diagonal), and radiative Weyl
    (antisymmetric off-diagonal) parts.
    $\mathrm{Tr}(\rho_{\mathrm{Weyl}}) = 0$ exactly.
    Two graviton polarisation modes emerge as degenerate eigenmodes.

  \item \textbf{Graviton amplitude theorem.}
    $\|W(\epsilon)\| = \sin(2\epsilon)/\sqrt{2}$, the exact GR analogue of the
    Born rule identity from Paper~VI.
    The graviton lifecycle --- born, peaked, vanished --- follows from the
    $\epsilon$-dependence, without postulating a graviton field.
\end{enumerate}

Together with Paper~VI, these results suggest that quantum mechanics and
general relativity are not two incompatible frameworks requiring unification.
They are two projections of the same compression principle onto different
physical degrees of freedom.
The wavefunction is the codec.
Measurement is decompression.
Particles are pixels.
Bosons and gravitons are compression residuals.
Spacetime geometry is what the codec looks like when applied to metric
configurations.

The identical formula $\sin(2\theta)/\sqrt{2}$ governing both boson amplitude (Paper VI) and graviton amplitude (this paper) is not a coincidence.
It is the unique consequence of the density matrix decomposition applied to any normalised configuration pair in superposition.
Quantum field theory and general relativity share this formula because they share the same codec.


% ============================================================
\section*{Supplementary Material}
\label{sec:supplementary}

Python experiments:

\begin{itemize}
  \item \texttt{ellipse\_theorem.py}.
  \item \texttt{graviton\_amplitude.py}.
  \item \texttt{kepler\_consistency.py}.
  \item \texttt{curvature\_decomposition.py}.
  \item \texttt{planck\_consistency.py}.
  \item \texttt{polarisation\_modes.py}.
\end{itemize}

\end{document}
